%%%%%%%%%%%%%%%%%%%%%%%%%%%%%%%%%%%%%%%%%%%%%%%%%%%%%%%%%%%%%%%%%%%%%%%%%%%%%%%%%%%%
% WASAFF CONSULTING - PREÁMBULO MAESTRO PARA PUBLICACIÓN TÉCNICA
% Versión: 2.0
% Propósito: Definir el entorno completo para la compilación de la monografía.
%%%%%%%%%%%%%%%%%%%%%%%%%%%%%%%%%%%%%%%%%%%%%%%%%%%%%%%%%%%%%%%%%%%%%%%%%%%%%%%%%%%%

% --- PAQUETES FUNDAMENTALES DE ESTRUCTURA Y LENGUAJE ---
\usepackage[utf8]{inputenc}
\usepackage[spanish,es-nosectiondot,es-lcroman]{babel}
\usepackage{geometry}
\usepackage{fancyhdr}
\usepackage{titlesec}
\usepackage{setspace}
\usepackage{parskip}

% --- PAQUETES MATEMÁTICOS Y DE SÍMBOLOS ---
\usepackage{amsmath}
\usepackage{amssymb}
\usepackage{bm}

% --- PAQUETES PARA ELEMENTOS GRÁFICOS Y VISUALES ---
\usepackage{graphicx}
\usepackage{xcolor}
\usepackage{tikz}
\usetikzlibrary{shapes.geometric, arrows.meta, positioning, babel, calc, decorations.pathmorphing}
\usepackage[most]{tcolorbox}
\tcbuselibrary{skins,listings,breakable}

% --- PAQUETES PARA TABLAS, CÓDIGO Y REFERENCIAS ---
\usepackage{booktabs}
\usepackage{listings}
\usepackage{hyperref}
\usepackage{microtype}

% --- CONFIGURACIÓN GLOBAL ---
\newcommand{\consultingAuthor}{Felipe Wasaff}
\newcommand{\consultingCompany}{Wasaff Consulting}
\geometry{a4paper, left=2.5cm, right=2.5cm, top=2.5cm, bottom=2.5cm}
\setlength{\headheight}{15pt}
\definecolor{WasaffNavy}{HTML}{001F3F}
\definecolor{WasaffCopper}{HTML}{B87333}
\definecolor{WasaffLightGrey}{HTML}{F0F0F0}
\definecolor{WasaffExampleBlue}{HTML}{EAF2F8}
\definecolor{WasaffNoteYellow}{HTML}{FFF9E5}
\definecolor{WasaffChallengeGreen}{HTML}{E8F8F5}
\definecolor{CodeBackground}{HTML}{FAFAFA}

% --- PORTADA ---
\title{\sffamily
    {\huge\bfseries\color{WasaffNavy} La Simulación Numérica como Herramienta de Ingeniería:}\\[0.5em]
    {\Large\color{black} Un Enfoque desde los Primeros Principios con FEniCS}
}
\author{\large Felipe Wasaff \\ \Large\textit{\color{WasaffCopper}\consultingCompany}}
\date{\large Edición \the\year}

% --- ESTILO DE CAPÍTULOS Y SECCIONES ---
\titleformat{\part}[display]
  {\normalfont\sffamily\huge\bfseries\color{WasaffCopper}}
  {\partname\ \thepart}
  {20pt}
  {\Huge\color{WasaffNavy}}
  [\vspace{1.5em}\noindent\color{black!40}\rule{\textwidth}{0.4pt}]

\titleformat{\chapter}[display]
  {\normalfont\sffamily\huge\bfseries\color{WasaffNavy}}
  {\fontsize{60}{60}\selectfont\color{WasaffCopper!30}\thechapter}
  {0pt}
  {\Huge}
\titlespacing*{\chapter}{0pt}{-40pt}{40pt}

\titleformat{\section}{\normalfont\sffamily\Large\bfseries\color{WasaffNavy}}{\thesection}{1em}{}
\titleformat{\subsection}{\normalfont\sffamily\large\bfseries\color{WasaffCopper}}{\thesubsection}{1em}{}

% --- ESTILO DE ENCABEZADO Y PIE DE PÁGINA ---
\pagestyle{fancy}
\fancyhf{}
\fancyhead[LE,RO]{\color{gray}\sffamily\thepage}
\fancyhead[RE]{\color{gray}\sffamily\nouppercase{\leftmark}}
\fancyhead[LO]{\color{gray}\sffamily\nouppercase{\rightmark}}
\renewcommand{\headrulewidth}{0.4pt}
\renewcommand{\footrulewidth}{0pt}
\fancypagestyle{plain}{
  \fancyhf{}%
  \renewcommand{\headrulewidth}{0pt}%
}

% --- HYPERREF Y METADATOS ---
\hypersetup{
    colorlinks=true, linkcolor=WasaffNavy, urlcolor=WasaffCopper,
    pdftitle={Simulación Numérica con FEniCS - Wasaff Consulting},
    pdfauthor={\consultingAuthor, \consultingCompany}
}

% --- ENTORNOS PERSONALIZADOS ---
\newtcolorbox{infobox}{breakable, enhanced, colback=WasaffLightGrey, colframe=WasaffNavy, fonttitle=\bfseries\sffamily, title=\MakeUppercase{Contexto Industrial}}
\newtcolorbox{keybox}[1]{breakable, enhanced, colback=#1!5, colframe=#1, fonttitle=\bfseries\sffamily\small, title=\MakeUppercase{#1}}
\newtcolorbox{notatecnica}{breakable, enhanced, colback=WasaffNoteYellow, colframe=WasaffCopper, fonttitle=\bfseries\sffamily, title=\MakeUppercase{Discusión Técnica}}
\newtcolorbox{challengebox}{breakable, enhanced, colback=WasaffChallengeGreen, colframe=green!40!black, fonttitle=\bfseries\sffamily, title=\MakeUppercase{Desafío de Certificación}}
\lstdefinestyle{pythonstyle}{backgroundcolor=\color{CodeBackground}, commentstyle=\color{green!40!black}, keywordstyle=\color{blue}, numberstyle=\tiny\color{gray}, stringstyle=\color{purple}, basicstyle=\ttfamily\footnotesize, breakatwhitespace=false, breaklines=true, captionpos=b, keepspaces=true, numbers=left, numbersep=5pt, showspaces=false, showstringspaces=false, showtabs=false, tabsize=2, language=Python, literate={á}{{\'a}}1 {é}{{\'e}}1 {í}{{\'i}}1 {ó}{{\'o}}1 {ú}{{\'u}}1 {Á}{{\'A}}1 {É}{{\'E}}1 {Í}{{\'I}}1 {Ó}{{\'O}}1 {Ú}{{\'U}}1 {ñ}{{\~n}}1 {Ñ}{{\~N}}1}